\documentclass[letterpaper]{article}

\usepackage[letterpaper,margin=1in]{geometry}

% Compile from the repository root so this path resolves to the pinned WSP
% submodule.
\input{wsp/testing/test-case-library.tex}

\begin{document}
\def\TCID{TC-0011}
\def\TCTitle{Run-time error messages}
\def\TCPurpose{Verify the messages of the BASIC run-time errors, the numbers
BASIC does not define, and the message an error raised by the runtime
carries.}
\def\TCPriority{\TCPriorityMedium}

\def\TCRequirementRef{REQ-0010}
\def\TCDesignRef{Not applicable.}
\def\TCFeatureRef{\texttt{include/wfc/brun71/error.hpp};
\texttt{tests/error\_test.cpp}.}
\def\TCTestTechnique{Sampling of the error table, including the messages
that differ from Visual Basic's; boundary values (0, 1, 77, a large
number).}

\def\TCPreconditions{The library is built and the test program
\texttt{brun71\_error\_test} exists.}
\def\TCEnvironment{Any supported target. CTest runs the program with a 60
second timeout.}
\def\TCAssumptions{The messages of the PDS 7.1 list of run-time error codes
(REQ-0010).}

\def\TCInputData{Error numbers 0, 1, 3, 5, 8, 11, 13, 16, 52, 62, 75, 76, 77
and 1000.}
\def\TCInitialState{None.}

\def\TCProcedure{\begin{enumerate}
  \item Run \texttt{brun71\_error\_test} (CTest \texttt{TC-0011-error}, in
    WFC \texttt{wfc-brun71-TC-0011-error}).
  \item Compare each message with the expected one, and check that the
    numbers BASIC does not define have none.
  \item Check the number and the message of a raised error.
\end{enumerate}}
\def\TCExpectedResult{The program exits with status 0 and prints
\texttt{errors: ok}; no \texttt{FAIL} line is written.}
\def\TCPostConditions{None; the test writes no files.}

\TCTemplate

\end{document}
